\section{Capsule two: keeping infinitely many indices}\label{ch17:sec:pigeonhole}
A sequence $x_0,x_1,x_2,\ldots$ has one term for each index $n\in\NN$, so it has infinitely many indices even when its values repeat: $x_n=(-1)^n$ takes only the values $1$ and $-1$, yet the value $1$ occurs at the infinitely many indices $0,2,4,\ldots$. Counting points instead would break the assistant's plan for a constant sequence: no ball contains infinitely many different terms, yet a constant sequence converges at once. So the proof keeps track of sets of \emph{indices}: a ball contains infinitely many terms of a sequence if it contains $x_n$ for infinitely many indices $n$, whether or not these terms are different points.

Recall from Section~\ref{ch03:sec:sets} that a set is \emph{finite} if its members can be written down in a list that ends, and \emph{infinite} otherwise.

\par\noindent\begin{minipage}{\linewidth}
\begin{lemma}[Infinite pigeonhole principle]\label{ch17:lem:pigeonhole}
Let $S$ be an infinite set, and let $A_1,\ldots,A_k$ be finitely many sets whose union contains $S$. Then at least one of the sets $S\cap A_1,\ldots,S\cap A_k$ is infinite.
\end{lemma}
\end{minipage}\par

\begin{proof}
Every member of $S$ lies in at least one of the sets $A_i$, so
\[
 S=(S\cap A_1)\cup(S\cap A_2)\cup\cdots\cup(S\cap A_k).
\]
Suppose, for contradiction, that each of the sets $S\cap A_i$ is finite. Write down the members of $S\cap A_1$, then those of $S\cap A_2$, and so on up to $S\cap A_k$. The result is a list that ends, and every member of $S$ appears in it at least once. Crossing out repetitions, which occur for members lying in several of the sets $A_i$, leaves a list of the members of $S$ that still ends. So $S$ is finite, a contradiction.
\end{proof}

This is the pigeonhole principle of Proposition~\ref{ch07:prop:pigeonhole} for infinitely many objects: if the members of $S$ are placed in finitely many boxes $A_1,\ldots,A_k$, some box holds infinitely many of them. Here, though, the boxes may overlap; an object in two boxes simply appears twice in the list.

For example, take $x_n=(-1)^n$ and the balls $B(1,1/2)$ and $B(-1,1/2)$ on the real line. The indices $n$ with $x_n\in B(1,1/2)$ form the set $A_1$ of even numbers, and those with $x_n\in B(-1,1/2)$ form the set $A_2$ of odd numbers. Their union is $\NN$, and the lemma with $S=\NN$ promises that at least one of them is infinite, without saying which. Here both are.

In the proof of the theorem, $S$ will be the current infinite set of indices and $A_i$ the set of indices $n$ for which $x_n$ lies in the $i$th ball of a finite cover; the lemma then supplies one ball containing $x_n$ for infinitely many $n\in S$.
